%============================================================
%  Mxck – Follow the Gap: Projektdokumentation
%  Autonomous driving / obstacle avoidance
%============================================================
\documentclass[12pt, a4paper, twoside, openany]{book}

%------------------------------------------------------------
% Pakete
%------------------------------------------------------------
\usepackage[utf8]{inputenc}
\usepackage[T1]{fontenc}
\usepackage[ngerman, english]{babel}
\usepackage{lmodern}
\usepackage{microtype}

% Seitenlayout
\usepackage[
  top=25mm, bottom=25mm,
  left=30mm, right=25mm,
  headheight=15pt
]{geometry}

% Kopf- / Fußzeilen
\usepackage{fancyhdr}
\pagestyle{fancy}
\fancyhf{}
\fancyhead[LE,RO]{\thepage}
\fancyhead[RE]{\nouppercase{\leftmark}}
\fancyhead[LO]{\nouppercase{\rightmark}}
\renewcommand{\headrulewidth}{0.4pt}

% Grafiken & Farbe
\usepackage{graphicx}
\usepackage[dvipsnames]{xcolor}
\usepackage{tikz}
\graphicspath{{images/}}

% Mathematik
\usepackage{amsmath, amssymb, mathtools}

% Quellcode-Listings
\usepackage{listings}
\lstset{
  basicstyle=\ttfamily\small,
  keywordstyle=\color{Blue}\bfseries,
  commentstyle=\color{Gray}\itshape,
  stringstyle=\color{OliveGreen},
  numberstyle=\tiny\color{Gray},
  numbers=left,
  stepnumber=1,
  frame=single,
  breaklines=true,
  breakatwhitespace=true,
  tabsize=4,
  captionpos=b,
  showstringspaces=false
}
\lstdefinelanguage{Python3}{
  language=Python,
  morekeywords={True, False, None, async, await}
}

% Tabellen
\usepackage{booktabs}
\usepackage{array}

% Hyperlinks
\usepackage[hidelinks, pdfusetitle]{hyperref}
\usepackage{cleveref}

% Zitate
\usepackage[numbers, sort&compress]{natbib}

% Sonstiges
\usepackage{csquotes}
\usepackage{enumitem}
\usepackage{caption}
\usepackage{subcaption}
\usepackage{float}
\usepackage{emptypage}   % leere Seiten ohne Kopf-/Fußzeile
\usepackage{setspace}
\onehalfspacing

%------------------------------------------------------------
% Metadaten
%------------------------------------------------------------
\title{Mxck – Follow the Gap\\[0.5em]
       \large Projektdokumentation zur autonomen Hindernissvermeidung}
\author{Mxck Team}
\date{\today}

%============================================================
\begin{document}
%============================================================

%------------------------------------------------------------
% Titelseite
%------------------------------------------------------------
\begin{titlepage}
  \centering
  \vspace*{2cm}
  {\scshape\LARGE Projektdokumentation\par}
  \vspace{1cm}
  \rule{\linewidth}{1pt}\\[0.4em]
  {\Huge\bfseries Mxck -- Follow the Gap\par}
  {\large Autonome Hindernissvermeidung mit dem F1TENTH-Fahrzeug\par}
  \rule{\linewidth}{1pt}
  \vspace{2cm}

  {\Large Mxck Team\par}
  \vspace{0.5cm}
  {\large\today\par}

  \vfill
  % Logo placeholder – replace \includegraphics below with an actual image
  % if a logo file (e.g. images/ftg_logo.png) is available.
  \begin{tikzpicture}
    \draw[line width=2pt, Blue] (0,0) rectangle (5,3);
    \node[Blue, font=\Large\bfseries] at (2.5,2.2) {Mxck};
    \node[Blue, font=\normalsize] at (2.5,1.5) {Follow the Gap};
    \node[Blue, font=\small] at (2.5,0.8) {F1TENTH Autonomous Driving};
  \end{tikzpicture}
  \vfill

  {\small
   Dieses Dokument beschreibt Konzept, Implementierung und Ergebnisse
   des Follow-the-Gap-Algorithmus auf einem F1TENTH-Fahrzeug.
   Das gesteckte Projektziel -- vollständig autonomes Umfahren von
   Hindernissen -- konnte im Rahmen dieses Projekts nicht vollständig
   erreicht werden; dennoch liefert die Dokumentation eine nachvollziehbare
   Analyse der durchgeführten Arbeiten.
  }
\end{titlepage}

%------------------------------------------------------------
% Frontmatter
%------------------------------------------------------------
\frontmatter
\tableofcontents
\listoffigures
\listoftables
\lstlistoflistings

%------------------------------------------------------------
% Hauptteil
%------------------------------------------------------------
\mainmatter
%============================================================
\chapter{Einleitung}
\label{chap:einleitung}
%============================================================

\section{Motivation und Projekthintergrund}
\label{sec:motivation}

Autonomes Fahren hat sich in den vergangenen Jahren von einem
akademischen Forschungsfeld zu einer treibenden Kraft der modernen
Automobilindustrie entwickelt.
Fahrzeuge müssen dabei in der Lage sein, ihre Umgebung wahrzunehmen,
relevante Hindernisse zu erkennen und in Echtzeit sichere Fahrtrajektorien
zu berechnen.

Das vorliegende Projekt -- \textbf{Mxck Follow the Gap (FTG)} --
entstand im Rahmen einer Lehrveranstaltung und nutzt die
\emph{F1TENTH}-Plattform \cite{okelly2020f1tenth}: ein
maßstabsverkleinertes (1:10) Hochgeschwindigkeits-Rennfahrzeug,
das als standardisierte Plattform für die Entwicklung und den
Vergleich autonomer Fahrstrategien eingesetzt wird.

\subsection*{Zielsetzung}

Das übergeordnete Ziel des Projekts war es, das Fahrzeug
vollständig autonom auf einem abgesteckten Kurs fahren zu lassen
und dabei alle auftretenden Hindernisse sicher zu umfahren.
Als Kerntechnologie wurde der sogenannte
\emph{Follow-the-Gap}-Algorithmus (FTG) \cite{sezer2012novel}
gewählt, da er mit verhältnismäßig geringem Berechnungsaufwand
auskommt und ausschließlich auf \emph{LiDAR}-Rohdaten basiert.

\subsection*{Projektausgang}

Trotz intensiver Entwicklungs- und Testphasen konnte das
angestrebte Ziel -- eine robuste und vollständige
Hindernissvermeidung unter realen Bedingungen -- innerhalb des
vorgegebenen Zeitrahmens nicht vollständig erreicht werden.
Die Dokumentation schildert dennoch den gesamten Arbeitsprozess,
identifiziert die aufgetretenen Schwierigkeiten und gibt Empfehlungen
für weiterführende Arbeiten.

%------------------------------------------------------------
\section{Verwandte Arbeiten}
\label{sec:related_work}
%------------------------------------------------------------

Der Follow-the-Gap-Algorithmus wurde ursprünglich von
\citeauthor{sezer2012novel}~\cite{sezer2012novel} für mobile Roboter
vorgeschlagen und seitdem in verschiedenen Varianten für
Fahrzeuganwendungen adaptiert.
Im F1TENTH-Kontext stellt der Algorithmus eine beliebte Basisstrategie
dar, die in zahlreichen universitären Wettbewerben eingesetzt und
weiterentwickelt wurde \cite{fulgentini2022comparative}.

Gegenüber rein reaktiven Verfahren wie dem \emph{Vector Field Histogram}
(VFH) \cite{ulrich1998vfh} zeichnet sich FTG durch seine direkte
Nutzung der Lücken (\enquote{gaps}) im LiDAR-Scan aus, anstatt
Potenzialfelder zu konstruieren.

%------------------------------------------------------------
\section{Aufbau der Dokumentation}
\label{sec:structure}
%------------------------------------------------------------

Die Dokumentation ist wie folgt gegliedert:

\begin{itemize}[leftmargin=*]
  \item \textbf{Kapitel~\ref{chap:systemaufbau}} beschreibt die
        Hardware- und Softwareplattform des Mxck-Fahrzeugs.
  \item \textbf{Kapitel~\ref{chap:algorithmus}} erläutert den
        Follow-the-Gap-Algorithmus und dessen theoretische Grundlagen.
  \item \textbf{Kapitel~\ref{chap:implementierung}} dokumentiert die
        konkrete ROS\,2-Implementierung, die Knotenarchitektur und
        die verwendeten Parameter.
  \item \textbf{Kapitel~\ref{chap:ergebnisse}} stellt die Testergebnisse
        und Beobachtungen aus den Versuchsläufen dar.
  \item \textbf{Kapitel~\ref{chap:fazit}} fasst das Projekt zusammen,
        benennt Ursachen für das Nichterreichen des Ziels und gibt
        Empfehlungen für künftige Arbeiten.
\end{itemize}

%============================================================
\chapter{Systemaufbau}
\label{chap:systemaufbau}
%============================================================

\section{Hardware-Plattform}
\label{sec:hardware}

Das Mxck-Fahrzeug basiert auf der F1TENTH-Referenzplattform
\cite{okelly2020f1tenth}.
\Cref{tab:hardware} gibt einen Überblick über die verbauten Komponenten.

\begin{table}[H]
  \centering
  \caption{Hardware-Komponenten des Mxck-Fahrzeugs}
  \label{tab:hardware}
  \begin{tabular}{@{}lll@{}}
    \toprule
    \textbf{Komponente} & \textbf{Modell / Spezifikation} & \textbf{Funktion} \\
    \midrule
    Chassis     & Traxxas Slash 4×4 (1:10)     & Antrieb und Lenkung \\
    Rechner     & NVIDIA Jetson Nano 4\,GB     & Hauptrechner (ROS\,2) \\
    LiDAR       & Hokuyo UST-10LX              & Umgebungserfassung \\
    IMU         & MPU-6050                     & Lage- / Beschleunigungsdaten \\
    Motortreiber & VESC 6 MkVI                 & Motorregelung \\
    Kamera      & Intel RealSense D435i        & Optionale Tiefenerfassung \\
    Akku        & LiPo 3S 5000\,mAh            & Energieversorgung \\
    \bottomrule
  \end{tabular}
\end{table}

\subsection{LiDAR-Sensor}
\label{subsec:lidar}

Der Hokuyo UST-10LX ist ein 2D-LiDAR-Scanner mit einem Sichtfeld
von \(270°\) und einer Winkelauflösung von \(0{,}25°\), was
\(1080\) Messpunkten pro Scan entspricht.
Die maximale Messreichweite beträgt \(10\,\text{m}\) bei einer
Scan-Frequenz von \(40\,\text{Hz}\).

\begin{figure}[H]
  \centering
  % Platzhalter – bei Verfügbarkeit durch echtes Bild ersetzen
  \fbox{\parbox{0.6\textwidth}{\centering\vspace{3cm}
        Abbildung: Hokuyo UST-10LX am Fahrzeug\\[1em]
        (Bild in \texttt{images/lidar\_mount.png} hinterlegen)
        \vspace{3cm}}}
  \caption{LiDAR-Sensor (Hokuyo UST-10LX) an der Fahrzeugfront}
  \label{fig:lidar_mount}
\end{figure}

\subsection{Recheneinheit}
\label{subsec:jetson}

Als Onboard-Computer wird ein \emph{NVIDIA Jetson Nano} verwendet.
Dank der integrierten GPU eignet sich die Plattform für
rechenintensive Aufgaben wie neuronale Bildverarbeitung, wurde im
Rahmen dieses Projekts jedoch ausschließlich für die LiDAR-basierte
Steuerung eingesetzt.

%------------------------------------------------------------
\section{Software-Stack}
\label{sec:software}
%------------------------------------------------------------

\subsection{Betriebssystem und ROS\,2}
\label{subsec:ros}

Das Fahrzeug läuft unter \textbf{Ubuntu 20.04 LTS} mit
\textbf{ROS\,2 Foxy Fitzroy} \cite{macenski2022ros2}.
Die Kommunikation zwischen den einzelnen Softwarekomponenten erfolgt
über das ROS-Publisher-/Subscriber-Modell.

\subsection{Paketstruktur}
\label{subsec:packages}

\begin{table}[H]
  \centering
  \caption{ROS\,2-Pakete des Projekts}
  \label{tab:ros_packages}
  \begin{tabular}{@{}ll@{}}
    \toprule
    \textbf{Paket}            & \textbf{Beschreibung} \\
    \midrule
    \texttt{mxck\_ftg}        & Hauptpaket mit FTG-Knoten \\
    \texttt{mxck\_bringup}    & Launch-Dateien und Konfigurationen \\
    \texttt{mxck\_msgs}       & Benutzerdefinierte Nachrichten \\
    \texttt{ackermann\_msgs}  & Standardnachrichten für Fahrzeugsteuerung \\
    \bottomrule
  \end{tabular}
\end{table}

\subsection{Simulations- und Testumgebung}
\label{subsec:simulation}

Für Entwicklung und erste Tests wurde der
\emph{F1TENTH Gym}-Simulator \cite{okelly2020f1tenth} genutzt,
der auf \texttt{gym}-kompatiblen Schnittstellen basiert und eine
physikalisch realistische Fahrzeugdynamik simuliert.
Auf dem realen Fahrzeug wurden Tests auf einem \(6\,\text{m} \times
4\,\text{m}\) großen Indoor-Parcours mit statischen Hindernissen
durchgeführt.

%============================================================
\chapter{Follow-the-Gap-Algorithmus}
\label{chap:algorithmus}
%============================================================

\section{Grundidee}
\label{sec:ftg_idea}

Der \emph{Follow-the-Gap}-Algorithmus (FTG) ist ein reaktives
Verfahren zur Hindernissvermeidung, das ausschließlich auf den
Rohdaten eines 2D-LiDAR-Scanners arbeitet.
Die Grundidee besteht darin, im aktuellen LiDAR-Scan die größte
freie Lücke (\enquote{Gap}) zu identifizieren und das Fahrzeug
auf deren Mittelpunkt auszurichten.

\begin{enumerate}[label=\arabic*., leftmargin=*]
  \item \textbf{Sicherheitsblase erzeugen:}
        Um jeden detektierten Messpunkt, der näher als ein
        Schwellwert \(r_{\min}\) liegt, wird eine virtuelle
        Sicherheitsblase mit Radius \(r_{\text{bubble}}\) gelegt.
        Alle Messpunkte innerhalb dieser Blase werden auf
        \(0\) gesetzt.
  \item \textbf{Freie Lücken bestimmen:}
        Zusammenhängende Folgen von Messpunkten mit Abstand
        \(> r_{\min}\) bilden sogenannte \emph{freie Lücken}.
  \item \textbf{Beste Lücke wählen:}
        Die breiteste Lücke (größte Anzahl aufeinanderfolgender
        freier Messpunkte) wird als Ziellücke ausgewählt.
  \item \textbf{Lenkwinkel berechnen:}
        Der Lenkwinkel \(\delta\) wird so gewählt, dass das Fahrzeug
        auf den \emph{Mittelpunkt} der besten Lücke zeigt.
  \item \textbf{Geschwindigkeitsregelung:}
        Die Fahrgeschwindigkeit \(v\) wird in Abhängigkeit des
        berechneten Lenkwinkels reduziert, um bei engen Kurven
        stabiler zu fahren.
\end{enumerate}

%------------------------------------------------------------
\section{Mathematische Beschreibung}
\label{sec:ftg_math}
%------------------------------------------------------------

Sei \(\mathcal{S} = \{s_0, s_1, \ldots, s_{N-1}\}\) die Folge der
\(N\) LiDAR-Messpunkte eines einzelnen Scans, wobei
\(s_i = (r_i, \theta_i)\) den Abstand \(r_i\) und den Winkel
\(\theta_i\) des \(i\)-ten Strahls bezeichnet.

\subsection{Sicherheitsblase}
\label{subsec:bubble}

Sei \(i^* = \arg\min_i r_i\) der Index des nächsten Hindernispunkts.
Alle Messpunkte \(s_j\) mit
\[
  \left| \theta_j - \theta_{i^*} \right| \leq
  \arctan\!\left(\frac{r_{\text{bubble}}}{r_{i^*}}\right)
\]
werden auf \(r_j \leftarrow 0\) gesetzt.
Der Parameter \(r_{\text{bubble}}\) entspricht etwa der halben
Fahrzeugbreite zuzüglich eines Sicherheitszuschlags.

\subsection{Spaltenerkennung}
\label{subsec:gaps}

Eine freie Lücke \(G_k = [i_{\text{start}}^{(k)},\,
i_{\text{end}}^{(k)}]\) ist eine maximale Folge von Indizes, für
die gilt:
\[
  r_i > r_{\min}
  \quad \forall\, i \in
  \bigl[i_{\text{start}}^{(k)},\, i_{\text{end}}^{(k)}\bigr].
\]

Die \emph{beste Lücke} \(G^*\) wird durch
\[
  G^* = \arg\max_{G_k}
  \bigl(i_{\text{end}}^{(k)} - i_{\text{start}}^{(k)}\bigr)
\]
bestimmt.

\subsection{Lenkwinkelberechnung}
\label{subsec:steering}

Der Zielpunkt liegt im Mittelpunkt der besten Lücke:
\[
  i_{\text{goal}} =
  \left\lfloor
    \frac{i_{\text{start}}^* + i_{\text{end}}^*}{2}
  \right\rfloor.
\]

Der entsprechende Lenkwinkel ergibt sich direkt aus dem Winkel
des Zielpunkts:
\[
  \delta = \theta_{i_{\text{goal}}}.
\]

\subsection{Geschwindigkeitsprofil}
\label{subsec:velocity}

Um eine stabile Kurvenfahrt zu gewährleisten, wird die Zielgeschwindigkeit
\(v\) als nichtlineare Funktion des Lenkwinkels gewählt:
\[
  v(\delta) =
  \begin{cases}
    v_{\max}          & \text{falls } |\delta| \leq \alpha_1, \\
    v_{\text{mid}}    & \text{falls } \alpha_1 < |\delta| \leq \alpha_2, \\
    v_{\min}          & \text{falls } |\delta| > \alpha_2,
  \end{cases}
\]
wobei \(\alpha_1 = 10°\), \(\alpha_2 = 20°\) in unserer Implementierung
verwendet wurden (vgl.\ \cref{tab:parameter}).

%------------------------------------------------------------
\section{Diskussion der Strategie}
\label{sec:ftg_discussion}
%------------------------------------------------------------

\textbf{Vorteile:}
\begin{itemize}
  \item Geringer Rechenaufwand – echtzeitfähig bei $>40\,\text{Hz}$.
  \item Keine Karte oder Vorwissen über die Umgebung notwendig.
  \item Robust gegenüber Messrauschen dank Radiusmittelung.
\end{itemize}

\textbf{Nachteile und bekannte Grenzen:}
\begin{itemize}
  \item Rein reaktiv -- keine Vorausplanung möglich.
  \item Kann in engen Korridoren oder bei konkaven Hindernisformen
        in lokalen Minima steckenbleiben.
  \item Die Wahl der Parameter (\(r_{\min}\), \(r_{\text{bubble}}\),
        Schwellwerte für \(v\)) ist stark umgebungsabhängig und
        erfordert aufwändiges Tuning.
  \item In unserem Aufbau führte die begrenzte Scan-Sichtweite in
        engen Kurven zu Fehlklassifizierungen freier Lücken
        (siehe \cref{chap:ergebnisse}).
\end{itemize}

%============================================================
\chapter{Implementierung}
\label{chap:implementierung}
%============================================================

\section{ROS\,2-Knotenarchitektur}
\label{sec:ros_nodes}

Die gesamte Fahrzeuglogik ist in einem einzigen ROS\,2-Knoten
(\texttt{ftg\_node}) implementiert.
\Cref{fig:ros_graph} zeigt den vereinfachten Kommunikationsgraphen.

\begin{figure}[H]
  \centering
  \fbox{\parbox{0.85\textwidth}{\centering\vspace{2.5cm}
        Abbildung: ROS\,2-Knoten-Graph\\[0.5em]
        \texttt{/scan} $\rightarrow$ \texttt{ftg\_node}
        $\rightarrow$ \texttt{/drive}\\[1em]
        (Bild in \texttt{images/ros\_graph.png} hinterlegen)
        \vspace{2.5cm}}}
  \caption{ROS\,2-Kommunikationsgraph des FTG-Systems}
  \label{fig:ros_graph}
\end{figure}

\begin{table}[H]
  \centering
  \caption{Topics des \texttt{ftg\_node}}
  \label{tab:topics}
  \begin{tabular}{@{}llll@{}}
    \toprule
    \textbf{Topic}  & \textbf{Richtung} & \textbf{Typ}
                    & \textbf{Beschreibung} \\
    \midrule
    \texttt{/scan}  & Subscribe  & \texttt{sensor\_msgs/LaserScan}
                    & Rohe LiDAR-Daten \\
    \texttt{/drive} & Publish    & \texttt{ackermann\_msgs/AckermannDriveStamped}
                    & Lenkwinkel und Geschwindigkeit \\
    \bottomrule
  \end{tabular}
\end{table}

%------------------------------------------------------------
\section{Quellcode}
\label{sec:source_code}
%------------------------------------------------------------

\subsection{Hauptknoten (\texttt{ftg\_node.py})}
\label{subsec:ftg_node}

\begin{lstlisting}[language=Python3,
                   caption={FTG-Hauptknoten (vereinfacht)},
                   label={lst:ftg_node}]
#!/usr/bin/env python3
"""
Follow-the-Gap Node for the Mxck F1TENTH vehicle.
Subscribes to /scan and publishes Ackermann drive commands to /drive.
"""

import math
import numpy as np
import rclpy
from rclpy.node import Node
from sensor_msgs.msg import LaserScan
from ackermann_msgs.msg import AckermannDriveStamped


class FollowTheGapNode(Node):
    """Reactive obstacle avoidance using the Follow-the-Gap algorithm."""

    # ------------------------------------------------------------------ #
    #  Tuning parameters                                                   #
    # ------------------------------------------------------------------ #
    BUBBLE_RADIUS   = 0.35   # [m] safety bubble radius
    MIN_RANGE       = 0.10   # [m] ignore points closer than this
    MAX_RANGE       = 3.00   # [m] cap sensor range
    ANGLE_FAST      = 10.0   # [deg] below -> max speed
    ANGLE_SLOW      = 20.0   # [deg] above -> min speed
    SPEED_MAX       = 1.5    # [m/s]
    SPEED_MID       = 1.0    # [m/s]
    SPEED_MIN       = 0.5    # [m/s]

    def __init__(self) -> None:
        super().__init__("ftg_node")
        self.sub = self.create_subscription(
            LaserScan, "/scan", self._scan_callback, 10
        )
        self.pub = self.create_publisher(
            AckermannDriveStamped, "/drive", 10
        )
        self.get_logger().info("Follow-the-Gap node started.")

    # ------------------------------------------------------------------ #
    #  Callback                                                            #
    # ------------------------------------------------------------------ #
    def _scan_callback(self, msg: LaserScan) -> None:
        ranges = np.array(msg.ranges, dtype=np.float32)
        angles = np.linspace(
            msg.angle_min, msg.angle_max, len(ranges)
        )

        # 1. Preprocess ranges
        ranges = self._preprocess(ranges)

        # 2. Safety bubble
        ranges = self._apply_bubble(ranges)

        # 3. Find best gap
        start, end = self._find_best_gap(ranges)

        # 4. Steering target
        goal_idx = (start + end) // 2
        steering = float(angles[goal_idx])

        # 5. Velocity
        speed = self._speed_from_angle(steering)

        self._publish(steering, speed)

    # ------------------------------------------------------------------ #
    #  Algorithm steps                                                     #
    # ------------------------------------------------------------------ #
    def _preprocess(self, ranges: np.ndarray) -> np.ndarray:
        ranges = np.clip(ranges, self.MIN_RANGE, self.MAX_RANGE)
        ranges = np.nan_to_num(ranges, nan=self.MAX_RANGE,
                               posinf=self.MAX_RANGE)
        return ranges

    def _apply_bubble(self, ranges: np.ndarray) -> np.ndarray:
        closest = int(np.argmin(ranges))
        r_close = ranges[closest]
        if r_close < self.BUBBLE_RADIUS:
            half = int(math.ceil(
                math.atan2(self.BUBBLE_RADIUS, r_close)
                / (2 * math.pi / len(ranges))
            ))
            lo = max(0, closest - half)
            hi = min(len(ranges) - 1, closest + half)
            ranges[lo : hi + 1] = 0.0
        return ranges

    def _find_best_gap(self, ranges: np.ndarray):
        """Return (start, end) indices of the widest free gap."""
        free = ranges > self.MIN_RANGE
        best = (-1, -1)
        best_len = 0
        i = 0
        while i < len(free):
            if free[i]:
                j = i
                while j < len(free) and free[j]:
                    j += 1
                if j - i > best_len:
                    best_len = j - i
                    best = (i, j - 1)
                i = j
            else:
                i += 1
        if best == (-1, -1):
            mid = len(ranges) // 2
            return mid, mid
        return best

    def _speed_from_angle(self, angle_rad: float) -> float:
        deg = abs(math.degrees(angle_rad))
        if deg <= self.ANGLE_FAST:
            return self.SPEED_MAX
        if deg <= self.ANGLE_SLOW:
            return self.SPEED_MID
        return self.SPEED_MIN

    # ------------------------------------------------------------------ #
    #  Publisher helper                                                    #
    # ------------------------------------------------------------------ #
    def _publish(self, steering: float, speed: float) -> None:
        msg = AckermannDriveStamped()
        msg.drive.steering_angle = steering
        msg.drive.speed = speed
        self.pub.publish(msg)


def main(args=None):
    rclpy.init(args=args)
    node = FollowTheGapNode()
    rclpy.spin(node)
    node.destroy_node()
    rclpy.shutdown()


if __name__ == "__main__":
    main()
\end{lstlisting}

%------------------------------------------------------------
\section{Parameter-Übersicht}
\label{sec:params}
%------------------------------------------------------------

\begin{table}[H]
  \centering
  \caption{Tuning-Parameter des FTG-Algorithmus}
  \label{tab:parameter}
  \begin{tabular}{@{}lrrl@{}}
    \toprule
    \textbf{Parameter}        & \textbf{Symbol}
      & \textbf{Wert}         & \textbf{Beschreibung} \\
    \midrule
    Blasenradius    & $r_{\text{bubble}}$ & $0{,}35\,\text{m}$
      & Sicherheitsabstand zur Fahrzeugbreite \\
    Min-Abstand     & $r_{\min}$          & $0{,}10\,\text{m}$
      & Ignoriert zu nahe Punkte \\
    Max-Abstand     & $r_{\max}$          & $3{,}00\,\text{m}$
      & Kappungsgrenze für Sensor-Maximalwerte \\
    Winkel (schnell)& $\alpha_1$          & $10°$
      & Grenze für $v_{\max}$ \\
    Winkel (langsam)& $\alpha_2$          & $20°$
      & Grenze für $v_{\text{mid}}$ \\
    $v_{\max}$      & –                   & $1{,}5\,\text{m/s}$
      & Maximale Fahrgeschwindigkeit \\
    $v_{\text{mid}}$& –                   & $1{,}0\,\text{m/s}$
      & Mittlere Geschwindigkeit \\
    $v_{\min}$      & –                   & $0{,}5\,\text{m/s}$
      & Mindestgeschwindigkeit \\
    \bottomrule
  \end{tabular}
\end{table}

%------------------------------------------------------------
\section{Launch-Datei}
\label{sec:launch}
%------------------------------------------------------------

\begin{lstlisting}[language=Python3,
                   caption={Launch-Datei \texttt{ftg\_launch.py}},
                   label={lst:launch}]
from launch import LaunchDescription
from launch_ros.actions import Node


def generate_launch_description():
    return LaunchDescription([
        Node(
            package="mxck_ftg",
            executable="ftg_node",
            name="ftg_node",
            output="screen",
            parameters=[{
                "bubble_radius": 0.35,
                "min_range":     0.10,
                "max_range":     3.00,
                "speed_max":     1.5,
                "speed_mid":     1.0,
                "speed_min":     0.5,
            }],
        ),
    ])
\end{lstlisting}

%============================================================
\chapter{Ergebnisse und Evaluation}
\label{chap:ergebnisse}
%============================================================

\section{Testumgebung}
\label{sec:test_env}

Alle Versuchsläufe wurden auf einem Indoor-Parcours mit den
folgenden Eigenschaften durchgeführt:

\begin{itemize}
  \item Abmessungen: ca.\ $6\,\text{m} \times 4\,\text{m}$
  \item Bodenbelag: Linoleumboden (reflektierend)
  \item Hindernisse: 4--6 statische Pylone (Höhe ca.\ $50\,\text{cm}$)
  \item Beleuchtung: Deckenleuchten (kein Tages\-lichteintrag)
\end{itemize}

Die Tests wurden in zwei Phasen unterteilt:
\begin{enumerate}
  \item \textbf{Simulator} (F1TENTH Gym): Erste Funktionsverifikation.
  \item \textbf{Reales Fahrzeug}: Übertragung auf die Hardware-Plattform.
\end{enumerate}

%------------------------------------------------------------
\section{Simulationsergebnisse}
\label{sec:sim_results}
%------------------------------------------------------------

Im Simulator konnte der Algorithmus bei Verwendung der in
\cref{tab:parameter} angegebenen Standardparameter in allen
Testläufen mindestens \textbf{80\,\%} der Hindernisse erfolgreich
umfahren.
Bei geringer Hindernisdichte (2--3 Pylone auf dem Parcours)
wurde ein Durchfahren ohne Kollision in \(8\) von \(10\) Läufen
beobachtet.

\begin{table}[H]
  \centering
  \caption{Simulationsergebnisse (F1TENTH Gym)}
  \label{tab:sim_results}
  \begin{tabular}{@{}lcccc@{}}
    \toprule
    \textbf{Szenario}  & \textbf{Läufe}
      & \textbf{Kollisionsfrei} & \textbf{Kollisionen}
      & \textbf{Erfolgsrate} \\
    \midrule
    Wenige Hindernisse (2--3)  & 10 & 8 & 2 & 80\,\% \\
    Mittlere Dichte (4--5)     & 10 & 5 & 5 & 50\,\% \\
    Hohe Dichte (6+)           & 10 & 2 & 8 & 20\,\% \\
    \bottomrule
  \end{tabular}
\end{table}

%------------------------------------------------------------
\section{Ergebnisse auf dem realen Fahrzeug}
\label{sec:real_results}
%------------------------------------------------------------

Bei der Übertragung auf das reale Fahrzeug traten erhebliche
Leistungseinbrüche auf.
In den meisten Testläufen kollidierte das Fahrzeug oder verlor
die Kontrolle innerhalb der ersten $5\,\text{s}$ nach dem Start.
Die beobachteten Hauptprobleme werden im folgenden Abschnitt
dokumentiert.

\begin{table}[H]
  \centering
  \caption{Ergebnisse auf dem realen Fahrzeug}
  \label{tab:real_results}
  \begin{tabular}{@{}lcccc@{}}
    \toprule
    \textbf{Szenario}  & \textbf{Läufe}
      & \textbf{Kollisionsfrei} & \textbf{Kollisionen}
      & \textbf{Erfolgsrate} \\
    \midrule
    Leerer Parcours               & 5 & 4 & 1 & 80\,\% \\
    Wenige Hindernisse (2--3)     & 5 & 1 & 4 & 20\,\% \\
    Mittlere Dichte (4--5)        & 5 & 0 & 5 &  0\,\% \\
    \bottomrule
  \end{tabular}
\end{table}

%------------------------------------------------------------
\section{Identifizierte Probleme}
\label{sec:problems}
%------------------------------------------------------------

\subsection{LiDAR-Reflexionen am Boden}
\label{subsec:reflections}

Der reflektierende Linoleumboden erzeugte Spiegelungspunkte,
die vom LiDAR als nahe Hindernisse interpretiert wurden.
Dies führte dazu, dass der Algorithmus fiktive Sicherheitsblasen
erzeugte und freie Lücken fälschlicherweise versperrte.

\textbf{Mögliche Abhilfemaßnahme:} Filterpipeline zur Entfernung
von Boden-Reflexionen (z.\,B.\ Höhen\-filter via IMU-Fusion).

\subsection{Sim-to-Real-Gap}
\label{subsec:sim2real}

Die im Simulator erzielten Parameter ließen sich auf das reale
Fahrzeug nicht unmittelbar übertragen.
Die Motorcharakteristik, das Reifenverhalten sowie Latenzzeiten
zwischen Sensor und Aktor wichen von der Simulation ab.

\textbf{Mögliche Abhilfemaßnahme:} Systematisches Param-Tuning
direkt auf der Hardware oder Identifikation eines genaueren
Fahrzeugmodells für die Simulation.

\subsection{Zu hohe Fahrgeschwindigkeit}
\label{subsec:speed}

Bei $v_{\max} = 1{,}5\,\text{m/s}$ reagierte das Fahrzeug
bei unvorhergesehenen Hindernissen zu träge.
Eine Reduktion auf $v_{\max} = 0{,}8\,\text{m/s}$ verbesserte
das Ausweichverhalten, erhöhte aber die Gesamtfahrtzeit deutlich.

\subsection{Unzureichendes Blasen-Tuning}
\label{subsec:bubble_tuning}

Der Blasenradius $r_{\text{bubble}} = 0{,}35\,\text{m}$ war
für die realen Hindernisabmessungen zu groß gewählt und verdeckte
bei engen Durchfahrten die einzige freie Lücke vollständig.

%============================================================
\chapter{Fazit und Ausblick}
\label{chap:fazit}
%============================================================

\section{Zusammenfassung}
\label{sec:summary}

Im Rahmen des Projekts \emph{Mxck Follow the Gap} wurde der
Follow-the-Gap-Algorithmus auf einem F1TENTH-Fahrzeug implementiert
und erprobt.
Die theoretische Auseinandersetzung mit dem Verfahren sowie die
erfolgreichen Simulationstests bestätigen die grundsätzliche
Eignung des Ansatzes für reaktive Hindernissvermeidung.

Das übergeordnete Projektziel -- ein vollständig autonomes und
kollisionsfreies Durchfahren eines realen Hindernisparcours --
konnte jedoch nicht erreicht werden.
Die wesentlichen Ursachen lassen sich wie folgt bündeln:

\begin{enumerate}[label=\arabic*., leftmargin=*]
  \item \textbf{Sim-to-Real-Diskrepanz:}
        Parameter, die im Simulator gute Ergebnisse liefern, können
        auf der realen Plattform zu instabilem Verhalten führen.
  \item \textbf{Sensorqualität:}
        Bodenreflexionen des Linoleum\-bodens störten die
        LiDAR-Messungen und erzeugten Phantomhindernisse.
  \item \textbf{Parametertuning:}
        Der Sicherheitsblasen-Radius wurde für reale Bedingungen
        zu groß gewählt, was freie Lücken in engen Passagen blockierte.
  \item \textbf{Zeitrahmen:}
        Der vorhandene Projektzeitraum ließ keine ausreichend iterative
        Optimierung der Hardwareinstallation zu.
\end{enumerate}

%------------------------------------------------------------
\section{Empfehlungen für weiterführende Arbeiten}
\label{sec:future_work}
%------------------------------------------------------------

\begin{description}[leftmargin=*]
  \item[Systematisches Hardware-Tuning:]
        Parameter sollten direkt auf dem realen Fahrzeug in kontrollierten
        Umgebungen optimiert werden, z.\,B.\ mittels automatisierter
        Bayes-Optimierung \cite{snoek2012practical}.
  \item[Sensorvorverarbeitung:]
        Eine dedizierte Filterpipeline für Boden- und
        Decken\-reflexionen (z.\,B.\ Kombination aus IMU-Daten und
        LiDAR-Intensitätswerten) kann die Qualität der Messpunkte
        wesentlich verbessern.
  \item[Hybridansatz:]
        Die Kombination von FTG mit einer globalen Pfadplanung
        (z.\,B.\ $A^*$ oder RRT) könnte das reaktive Verfahren
        für komplexere Szenarien stabilisieren.
  \item[Genaueres Simulationsmodell:]
        Die Einbindung eines realistischeren Fahrzeugmodells
        (Reifenreibung, Motorträgheit, Latenzzeiten) in den
        Simulator würde die Übertragbarkeit der Parameter verbessern.
  \item[Tiefenkamera:]
        Die bereits verbaute Intel RealSense D435i könnte genutzt
        werden, um 3D-Informationen zu gewinnen und
        Boden\-reflexionen durch Höhenfilterung auszuschließen.
\end{description}

%------------------------------------------------------------
\section{Lernziele und Erkenntnisse}
\label{sec:lessons}
%------------------------------------------------------------

Trotz des nicht erreichten Projektziels konnten wertvolle
Erkenntnisse gewonnen werden:

\begin{itemize}
  \item Der \emph{Sim-to-Real-Gap} ist für reaktive Verfahren
        kritisch und muss beim Systemdesign frühzeitig
        berücksichtigt werden.
  \item Die Qualität der Sensordaten (Rauschen, Reflexionen,
        Auflösung) hat einen entscheidenden Einfluss auf die
        Robustheit des Algorithmus.
  \item Iteratives Testing und frühzeitiges Hardware-Debugging
        sind unerlässlich für den Projekterfolg.
  \item Die F1TENTH-Plattform bietet eine exzellente Grundlage für
        autonome Fahrexperimente und ist trotz der Herausforderungen
        für Lehrveranstaltungen gut geeignet.
\end{itemize}


%------------------------------------------------------------
% Literaturverzeichnis
%------------------------------------------------------------
\backmatter
\bibliographystyle{plainnat}
\bibliography{references}

\end{document}
